\item Duality makes the middle form unimodular.  Its signature is two and
its positive and negative indices are three and one.  The rank-at-most-four integral classification
\cite[Theorem 2.2.1, p.~61]{ffPetersSterk} applied to these
positive and negative indices gives the odd form \(I_{3,1}\).
Since $H^2(M_c;\mathbb Z/2)=0$, its $w_2$ is zero.
For the complex tangent bundle,
$w_4=c_2\bmod2=(p_1/2)\bmod2$.
The Wu identity in degree four, with $w_1=w_2=0$, says that $w_4$ is
characteristic for the middle pairing.  Thus the vector $p_1/2$ is
characteristic and has square $2$.

